\Article{56}
Although the theory of quadratic residues and nonresidues in the arithmetic of complex numbers is contained within the theory of biquadratic residues itself, we shall present its principal theorems separately here before turning to the latter. For brevity, however, we shall discuss only the principal case, where the modulus is an odd complex prime number.

Let $m$ be such a modulus, and $p$ its norm. Evidently, any integer not divisible by $m$, a condition always to be understood here, can or cannot be congruent to a square modulo $m$, according as its index, with some primitive root taken as the base, is even or odd. In the former case that integer is called a quadratic residue of $m$, and in the latter a nonresidue. It follows that, among $p-1$ numbers forming a complete system of incongruent residues not divisible by $m$, one half belongs to the quadratic residues and the other half to the quadratic nonresidues. Any other number outside this system must be assigned the same
\SourceLeaf{137}%
character in this respect as the number in the system congruent to it.

It also follows that the product of two quadratic residues, as well as the product of two nonresidues, is a quadratic residue, whereas the product of a quadratic residue and a nonresidue is a nonresidue. In general, a product of any number of factors is a quadratic residue or a nonresidue according as the number of nonresidues among its factors is even or odd.

The following general criterion for distinguishing quadratic residues from nonresidues presents itself at once:

A number $k$ not divisible by the modulus is a quadratic residue or a nonresidue of it according as $k^{\frac12(p-1)}\equiv1$ or $k^{\frac12(p-1)}\equiv-1\pmod m$.

This theorem follows at once from the fact that, with any primitive root taken as the base, the index of the power $k^{\frac12(p-1)}$ becomes either $\equiv0$ or $\equiv\frac12(p-1)$, according as the index of $k$ is even or odd.

\Article{57}
For a given modulus, it is easy to divide a complete system of incongruent residues into two classes, quadratic residues and nonresidues. This also assigns all the remaining numbers to their respective classes. But a much deeper investigation is needed to find criteria distinguishing the moduli for which a given number is a quadratic residue from those for which it is a nonresidue.

As for the real units $+1$ and $-1$, these are actual squares in the arithmetic of complex numbers, and therefore quadratic residues for \emph{every} modulus. It follows just as easily from the criterion in the preceding article that $i$, and likewise $-i$, is a quadratic residue of every modulus whose norm $p$ has the form $8n+1$, but a nonresidue of every modulus whose norm has the form $8n+5$. Since it evidently makes no difference whether $m$ or one of its associates $im,-m,-im$ is taken as the modulus, we may suppose that the modulus is the primary one among its associates (Article~36, II). Thus, putting the modulus $=a+bi$, we may suppose $a$ odd and $b$ even. We then always have $aa\equiv1\pmod8$, while $bb$ is either $\equiv0$ or $\equiv4\pmod8$, according as $b$ is divisible by four or twice an odd number. Thus $+i$ and $-i$ are evidently quadratic residues of the modulus in the former case and nonresidues in the latter.

\SourceLeaf{138}%
\Article{58}
Since the character of a composite number, whether it is a quadratic residue or a nonresidue, depends on the characters of its factors, it evidently suffices to restrict the development of criteria distinguishing the moduli for which a given number $k$ is a quadratic residue from those for which it is a nonresidue to values of $k$ that are prime numbers and, moreover, primary among their associates. In this investigation, \emph{examining examples} immediately supplies very elegant theorems.

Let us begin with $1+i$, which is found to be a quadratic residue of the moduli
\[
\begin{gathered}
-1+2i,\ +3-2i,\ -5-2i,\ -1-6i,\ +5+4i,\ +5-4i,\\
-7,\ +7+2i,\ -5+6i,\quad\text{etc.}
\end{gathered}
\]
but a quadratic nonresidue of the following:
\[
\begin{gathered}
-1-2i,\ -3,\ +3+2i,\ +1+4i,\ +1-4i,\ -5+2i,\\
-1+6i,\ +7-2i,\ -5-6i,\ -3+8i,\ -3-8i,\\
+5+8i,\ +5-8i,\ +9+4i,\ +9-4i\quad\text{etc.}
\end{gathered}
\]

Examining this list carefully, in which we have always entered the primary one of each four associated moduli, we readily notice that all the moduli $a+bi$ in the first class satisfy $a+b$ $\equiv+1\pmod8$, while those in the second satisfy $a+b\equiv-3\pmod8$. If, in place of the primary modulus $m$, we take its associate $-m$, this criterion must evidently change so that $a+b\equiv-1$ for moduli in the first class, and $\equiv+3\pmod8$ for those in the second. Thus, provided our examination of examples has not misled us, in general, if $a+bi$ denotes a prime number in which $a$ is odd. $b$ even, $1+i$ is a quadratic residue or a quadratic nonresidue of it according as $a+b\equiv\pm1$ or $\equiv\pm3\pmod8$.

The same rule holds for $-1-i$ as for $1+i$. On the other hand, considering $1-i$ as the product of $-i$ and $1+i$, $1-i$ evidently has the same character as must be assigned to $1+i$ whenever $b$ is divisible by four, but the opposite character whenever $b$ is twice an odd number. It follows easily that $1-i$ is a quadratic residue of the prime number $a+bi$ whenever $a-b\equiv\pm1$, but a nonresidue whenever $a-b\equiv\pm3\pmod8$, always supposing $a$ odd and $b$ even.

This second proposition can also be deduced from the first with the help of a more general theorem, which we state as follows:

\SourceLeaf{139}%
In the theory of quadratic residues, the character of $\alpha+\beta i$ with respect to the modulus $a+bi$ is the same as that of $\alpha-\beta i$ with respect to the modulus $a-bi$.

The proof of this theorem follows from the fact that both moduli have the same norm $p$, and that whenever $(\alpha+\beta i)^{\frac12(p-1)}-1$ is divisible by $a+bi$, $(\alpha-\beta i)^{\frac12(p-1)}-1$ is also divisible by $a-bi$, while whenever $(\alpha+\beta i)^{\frac12(p-1)}+1$ is divisible by $a+bi$, $(\alpha-\beta i)^{\frac12(p-1)}+1$ must also be divisible by $a-bi$.

\Article{59}
Let us proceed to odd prime numbers.

We find that $-1+2i$ is a quadratic residue of the moduli $+3+2i$, $+1-4i$, $-5+2i$, $-5-2i$, $-1-6i$, $+7-2i$, $-3+8i$, $+5+8i$, $+5-8i$, $+9+4i$, etc.,

but a nonresidue of the moduli $-1-2i$, $-3$, $+3-2i$, $+1+4i$, $-1+6i$, $+5+4i$, $+5-4i$, $-7$, $+7+2i$, $-5+6i$, $-5-6i$, $-3-8i$, $+9-4i$, etc.

Reducing the moduli of the first class to their absolutely least residues modulo $-1+2i$, we find only $+1$ and $-1$, namely $+3+2i\equiv-1$, $+1-4i\equiv-1$, $-5+2i\equiv+1$, $-5-2i\equiv-1$, etc.

On the other hand, all moduli of the second class are found to be congruent modulo $-1+2i$ either to $+i$ or to $-i$.

But $+1,-1$ are themselves quadratic residues of the modulus $-1+2i$, while $+i$ and $-i$ are its nonresidues. Thus, as far as the examination of examples can be trusted, the following theorem emerges: $-1+2i$ is a quadratic residue or a quadratic nonresidue of the prime number $a+bi$ according as the latter is a quadratic residue or a quadratic nonresidue of $-1+2i$, provided that $a+bi$ is the primary one of its four associates, or rather, that $a$ is odd and $b$ even.

Analogous theorems concerning $+1-2i$, $-1-2i$, $+1+2i$ follow directly from this theorem.

\Article{60}
Examining examples similarly for $-3$ or $+3$, we find that each is a quadratic residue of the moduli $+3+2i$, $+3-2i$,
\SourceLeaf{140}%
$-1+6i$, $-1-6i$, $-7$, $-5+6i$, $-5-6i$, $-3+8i$, $-3-8i$, $+9+4i$, $+9-4i$, etc.,

but a nonresidue of $-1+2i$, $-1-2i$, $+1+4i$, $+1-4i$, $-5+2i$, $-5-2i$, $+5+4i$, $+5-4i$, $+7+2i$, $+7-2i$, $+5+8i$, $+5-8i$, etc.

\begingroup\rightskip=0pt plus 1fil\relax
The former are congruent modulo $3$ to one of the four numbers $+1,-1,+i,-i$, the latter to one of $+1+i,+1-i,-1+i,-1-i$. The former are precisely the quadratic residues of $3$, and the latter its nonresidues.\par\endgroup

Thus this examination of examples shows that the prime number $a+bi$, supposing $a$ odd and $b$ even, has the same relation to $-3$, and also to $+3$, as the latter has to it, in respect of whether one is a quadratic residue or a nonresidue of the other.

Extending a similar examination to other prime numbers, we find this most elegant law of reciprocity confirmed everywhere, and are led to the following fundamental theorem about quadratic residues in the arithmetic of complex numbers:

\emph{If $a+bi,A+Bi$ denote prime numbers such that $a,A$ are odd and $b,B$ even, either each is a quadratic residue of the other, or each is a nonresidue of the other.}

Yet despite the theorem's great simplicity, its proof is beset by great difficulties. We shall not dwell on these here, however, since the theorem itself is only a special case of a more general theorem that, as it were, exhausts the substance of the theory of biquadratic residues. Let us now turn to that theory.

